\begin{abstract}
Subtracting each participant's mean feature vector is a simple remedy for the between-participant variability that dominates cross-participant EEG decoding. We treat the step as an estimate of a participant offset from $m$ unlabeled reference observations and derive its consequences. The squared separation of two conditions after centering is $m/(m+1)$ times its asymptote, centering helps only above a break-even reference size $m^\ast$ fixed by the share of discriminant variance that belongs to participants, and in a complete repeated-measures panel leave-one-out centering is an exact rescaling of full-panel centering. On single trials of five public ERP CORE paradigms the law predicts accuracy at six reference sizes with a mean absolute error of 0.004, and the predicted break-even size agrees with the observed one in all five. Single-trial gains are small, from 0.000 to 0.036 balanced accuracy. On trial-averaged affective ERPs of 211 participants, where the participant share is large, centering on a three-condition panel raises three-class accuracy from 0.679 to 0.900, a two-condition panel gives 0.820, and a single observation gives none of the gain. Euclidean alignment recovers 0.834 on averages and nothing on single trials, which separates a participant gain from a participant offset. The results specify when subject centering is a valid description of deployment: the test participant must supply unlabeled reference observations, and their number sets the gain.
\end{abstract}

\begin{IEEEkeywords}
EEG decoding, cross-subject generalization, subject centering, transductive normalization, Euclidean alignment, event-related potentials, repeated measures.
\end{IEEEkeywords}
